% Source PDF: source/普通高考/2015/2015四川文.pdf, page 1.
% PDF embedded bitmap attempt: pdfimages -list returned no images; the flowchart is vector page content.
% Grid evidence: tmp/review_2015_sichuan_liberal/grids/page-grid100.png,
% tmp/review_2015_sichuan_liberal/grids/q6-grid50.png and q6-local-grid25.png.
% No-grid source crop: tmp/review_2015_sichuan_liberal/crops/q6_original_tight.png.
% Complete directed edges: start -> init(k=1) -> increment(k=k+1) -> test(k>4);
% test yes -> formula(S=sin(k*pi/6)) -> output -> finish; test no -> increment.
% The loop return joins the vertical entry above the increment node; no other edges exist.
% solid segments: all node borders and connector/loop segments; dashed segments: none.
% Labels: 是 is left of the true downward branch, 否 is right of the false loop branch;
% node text is centered with local zero indentation and no manual horizontal offset.
\begin{tikzpicture}[baseline,>=Stealth,font=\normalsize,line width=0.45pt,
  flowtext/.style={anchor=center,text centered,
    execute at begin node={\setlength{\parindent}{0pt}\noindent}},
  flowbox/.style={flowtext,draw,inner xsep=4pt,inner ysep=2.5pt}]
  \node[flowbox,rounded corners=5pt,text width=0.90cm,minimum width=0.90cm,minimum height=0.56cm] (start) at (0,0) {开始};
  \node[flowbox,text width=1.00cm,minimum width=1.00cm,minimum height=0.56cm] (init) at (0,-1.14) {$k=1$};
  \node[flowbox,text width=1.75cm,minimum width=1.75cm,minimum height=0.56cm] (inc) at (0,-2.27) {$k=k+1$};
  \node[flowbox,diamond,aspect=3.1,minimum width=2.70cm,minimum height=0.72cm] (test) at (0,-3.48) {$k>4?$};
  \node[flowbox,text width=2.60cm,minimum width=2.60cm,minimum height=1.15cm] (formula) at (0,-4.83) {$S=\sin\dfrac{k\pi}{6}$};
  \node[flowbox,trapezium,trapezium left angle=75,trapezium right angle=105,
        text width=1.50cm,minimum width=1.50cm,minimum height=0.58cm] (output) at (0,-6.12) {输出 $S$};
  \node[flowbox,rounded corners=5pt,text width=0.90cm,minimum width=0.90cm,minimum height=0.56cm] (finish) at (0,-7.25) {结束};

  \draw[->] (start.south) -- (init.north);
  % The loop return joins the incoming connector above k=k+1; place its
  % leftward arrowhead on the return segment before the merge, followed by
  % an unmarked tail into the single downward arrow.
  \coordinate (merge-level) at (0,-1.68);
  \coordinate (merge-tip) at (0.38,-1.68);
  \draw (init.south) -- (merge-level);
  \draw[->] (merge-level) -- (inc.north);
  \draw[->] (inc.south) -- (test.north);
  \draw[->] (test.south) -- (formula.north);
  \draw[->] (formula.south) -- (output.north);
  \draw[->] (output.south) -- (finish.north);
  \draw (test.east) -- (1.65,-3.48) -- (1.65,-1.68);
  \draw[->] (1.65,-1.68) -- (merge-tip);
  \draw (merge-tip) -- (merge-level);
  \node[flowtext,anchor=west,inner xsep=1pt,inner ysep=1pt] at (1.90,-3.12) {否};
  \node[flowtext,inner xsep=1pt,inner ysep=1pt] at (-0.42,-3.98) {是};
\end{tikzpicture}
